\scanpage{96}\mspage{40}
% L1-2
dans $\underline{\mathrm{Fibsc}}_{\underset{\uparrow\uparrow\nwarrow}{0}}\,S$, y-a-t-il une homotopie dans $\underline{\mathrm{Fibsc}}_{\underset{\uparrow\uparrow\uparrow}{0}}\,S$, ``naturelle'' i.e. fonctorielle en $\unsure{x|X}$,
% L3-5
\diagram{\begin{tikzcd}[ampersand replacement=\&,column sep=small]
 X\arrow[r]\&{}\arrow[r,"{\deleted{\unsure{\Phi}}\inserted{c}}",no head]\&
 \Phi(X)=\begin{gathered}\text{Catégorie \deleted{\unsure{co}}\unsure{fibrée|cofibrée}}\\
 \text{scindée de \unsure{fibre|fibres}}\\ \text{les }{}_{s\backslash}X\end{gathered}
\end{tikzcd}}
% L6-7
Mais un chemin qui soit formé de flèches
\[
 X=X_0\mathbin{-}X_1\mathbin{-}X_2\ldots\mathbin{-}X_n=\Phi(X)
\]
% L8-12
qui soient toutes \uline{cartésiennes} (et \unsure{de} plus, bien \unsure{sur|sûr}, dans $W$ !). Cela signifie donc qu'on \deleted{a des \unsure{c}} aurait des \deleted{\unsure{esp}} cat. fibrées scindées sur $S$, $X_1,\ldots,X_{n-1}$, dépendant \uline{fonctoriellement} de $X$, qui relient
% L13-14
\[
 X=X_0\qquad \text{à}\qquad \Phi(X)=X_n
\]
i.e.
\[
 X_s\qquad\text{à}\qquad {}_{s\backslash}X_{\deleted{\unsure{/s}}}\qquad ???
\]
% L15
Je ne vois rien à l'horizon
% archive sheet number: 95
